% GENERATED FILE. Edit canonical lessons, then run npm run generate:books.
% canonical-source-hash: fnv1a64:4f57b2efa3e735a9
% canonical-lessons: ML-C498-kantalkkatu, ML-C498-catuppu, ML-C498-pavilappurru, ML-C498-sravu, ML-C498-atinnukutuka

\chapter{Mangroves and Reefs}
\label{ch:ml-mangroves-and-reefs}
\hlchaptermodality{498}

\begin{chapteropening}
\textbf{By the end of this chapter:} \emph{I can say a mangrove forest, a swamp, a coral reef, a shark and to accumulate.}

Complete the last lesson of chapter 498: I can say a mangrove forest, a swamp, a coral reef, a shark and to accumulate.
\end{chapteropening}

\section[\texorpdfstring{\ml{കണ്ടൽക്കാട്} (kaṇṭalkkāṭŭ)}{കണ്ടൽക്കാട് (kaṇṭalkkāṭŭ)}]{\ml{കണ്ടൽക്കാട്} (kaṇṭalkkāṭŭ) — a mangrove forest}

\label{lesson:ML-C498-kantalkkatu}



\begin{quote}
Before the new one: say the Malayalam for low tide, then the Malayalam for a sea surge.
\end{quote}

\subsection*{You'll want to know: \ml{കണ്ടൽക്കാട്}}
\textbf{\ml{കണ്ടൽക്കാട്}} — \emph{kaṇṭalkkāṭŭ} — \textquotedblleft{}a mangrove forest\textquotedblright{}.

\begin{quote}
\ml{കണ്ടൽക്കാടുകൾ} \ml{തീരത്തെ} \ml{കടലിൽ} \ml{നിന്ന്} \ml{സംരക്ഷിക്കുന്നു}.
\end{quote}

\begin{quote}
\emph{kaṇṭalkkāṭukaḷ tīratte kaṭalil ninnŭ saṁrakṣikkunnu.}
\end{quote}

\begin{quote}
Mangrove forests protect the shore from the sea.
\end{quote}

\ml{കണ്ടൽ} is the mangrove tree; \ml{കാട്} is a forest.

\begin{cousinweb}
\textbf{\ml{കണ്ടൽക്കാട്}} (\emph{kaṇṭalkkāṭŭ}) is a native Dravidian compound of \ml{കണ്ടൽ}, a mangrove, and \ml{കാട്}, a forest (DEDR 1438). DEDR 1171 gives Tamil \emph{kaṇṭal} as a mangrove and lists the Malayalam word with the mangrove \emph{Rhizophora candel}; the same article also holds words for a bulbous root.
\end{cousinweb}

\begin{grammarlens}[title={What you've built}]
One more word for getting around.
\end{grammarlens}

\subsection*{Your turn}
\begin{itemize}
\raggedright
  \item \emph{Say it:} \emph{kaṇṭalkkāṭŭ}
  \item \emph{Say it:} \emph{kaṇṭalkkāṭŭ}, once more
  \item \emph{Say it:} \emph{kaṇṭalkkāṭukaḷ tīratte kaṭalil ninnŭ saṁrakṣikkunnu.} — Mangrove forests protect the shore from the sea.
  \item \emph{Recall:} say the Malayalam for low tide, then the Malayalam for a sea surge, then say \emph{kaṇṭalkkāṭŭ} again
\end{itemize}

\begin{tcolorbox}[breakable,colback=teal!4,colframe=teal!35!black,title={Before you move on}]
What does \ml{കണ്ടൽക്കാട്} mean? (\textquotedblleft{}A mangrove forest\textquotedblright{}.) Say it once more.
\end{tcolorbox}
\section[\texorpdfstring{\ml{ചതുപ്പ്} (catuppŭ)}{ചതുപ്പ് (catuppŭ)}]{\ml{ചതുപ്പ്} (catuppŭ) — a swamp; a marsh}

\label{lesson:ML-C498-catuppu}



\begin{quote}
Before the new one: say the Malayalam for a sea surge, then the Malayalam for a mangrove forest.
\end{quote}

\subsection*{You'll want to know: \ml{ചതുപ്പ്}}
\textbf{\ml{ചതുപ്പ്}} — \emph{catuppŭ} — \textquotedblleft{}a swamp; a marsh\textquotedblright{}.

\begin{quote}
\ml{ആ} \ml{ചതുപ്പിൽ} \ml{ധാരാളം} \ml{പക്ഷികൾ} \ml{ഉണ്ട്}.
\end{quote}

\begin{quote}
\emph{ā catuppil dhārāḷaṁ pakṣikaḷ uṇṭŭ.}
\end{quote}

\begin{quote}
There are many birds in that swamp.
\end{quote}

Wet land where your feet sink: \ml{ചതുപ്പുനിലം}.

\begin{cousinweb}
\textbf{\ml{ചതുപ്പ്}} (\emph{catuppŭ}) is native Dravidian, \textquotedblleft{}a marsh, boggy ground\textquotedblright{} (DEDR 2326).
\end{cousinweb}

\begin{grammarlens}[title={What you've built}]
One more word for getting around.
\end{grammarlens}

\subsection*{Your turn}
\begin{itemize}
\raggedright
  \item \emph{Say it:} \emph{catuppŭ}
  \item \emph{Say it:} \emph{catuppŭ}, once more
  \item \emph{Say it:} \emph{ā catuppil dhārāḷaṁ pakṣikaḷ uṇṭŭ.} — There are many birds in that swamp.
  \item \emph{Recall:} say the Malayalam for a sea surge, then the Malayalam for a mangrove forest, then say \emph{catuppŭ} again
\end{itemize}

\begin{tcolorbox}[breakable,colback=teal!4,colframe=teal!35!black,title={Before you move on}]
What does \ml{ചതുപ്പ്} mean? (\textquotedblleft{}A swamp; a marsh\textquotedblright{}.) Say it once more.
\end{tcolorbox}
\section[\texorpdfstring{\ml{പവിഴപ്പുറ്റ്} (paviḻappuṟṟŭ)}{പവിഴപ്പുറ്റ് (paviḻappuṟṟŭ)}]{\ml{പവിഴപ്പുറ്റ്} (paviḻappuṟṟŭ) — a coral reef}

\label{lesson:ML-C498-pavilappurru}



\begin{quote}
Before the new one: say the Malayalam for a mangrove forest, then the Malayalam for a swamp.
\end{quote}

\subsection*{You'll want to know: \ml{പവിഴപ്പുറ്റ്}}
\textbf{\ml{പവിഴപ്പുറ്റ്}} — \emph{paviḻappuṟṟŭ} — \textquotedblleft{}a coral reef\textquotedblright{}.

\begin{quote}
\ml{കടലിന്റെ} \ml{അടിയിൽ} \ml{സുന്ദരമായ} \ml{പവിഴപ്പുറ്റുകൾ} \ml{ഉണ്ട്}.
\end{quote}

\begin{quote}
\emph{kaṭalinṟe aṭiyil sundaramāya paviḻappuṟṟukaḷ uṇṭŭ.}
\end{quote}

\begin{quote}
There are beautiful coral reefs under the sea.
\end{quote}

\ml{പവിഴം}, coral, plus \ml{പുറ്റ്}, a mound: a mound of coral.

\begin{cousinweb}
\textbf{\ml{പവിഴപ്പുറ്റ്}} (\emph{paviḻappuṟṟŭ}) is a native Dravidian compound of \ml{പവിഴം}, coral (DEDR 3998), and \ml{പുറ്റ്}, a mound such as an anthill (DEDR 4335).
\end{cousinweb}

\begin{grammarlens}[title={What you've built}]
One more word for getting around.
\end{grammarlens}

\subsection*{Your turn}
\begin{itemize}
\raggedright
  \item \emph{Say it:} \emph{paviḻappuṟṟŭ}
  \item \emph{Say it:} \emph{paviḻappuṟṟŭ}, once more
  \item \emph{Say it:} \emph{kaṭalinṟe aṭiyil sundaramāya paviḻappuṟṟukaḷ uṇṭŭ.} — There are beautiful coral reefs under the sea.
  \item \emph{Recall:} say the Malayalam for a mangrove forest, then the Malayalam for a swamp, then say \emph{paviḻappuṟṟŭ} again
\end{itemize}

\begin{tcolorbox}[breakable,colback=teal!4,colframe=teal!35!black,title={Before you move on}]
What does \ml{പവിഴപ്പുറ്റ്} mean? (\textquotedblleft{}A coral reef\textquotedblright{}.) Say it once more.
\end{tcolorbox}
\section[\texorpdfstring{\ml{സ്രാവ്} (srāvŭ)}{സ്രാവ് (srāvŭ)}]{\ml{സ്രാവ്} (srāvŭ) — a shark}

\label{lesson:ML-C498-sravu}



\begin{quote}
Before the new one: say the Malayalam for a swamp, then the Malayalam for a coral reef.
\end{quote}

\subsection*{You'll want to know: \ml{സ്രാവ്}}
\textbf{\ml{സ്രാവ്}} — \emph{srāvŭ} — \textquotedblleft{}a shark\textquotedblright{}.

\begin{quote}
\ml{മീൻപിടുത്തക്കാരുടെ} \ml{വലയിൽ} \ml{ഒരു} \ml{സ്രാവ്} \ml{പെട്ടു}.
\end{quote}

\begin{quote}
\emph{mīnpiṭuttakkāruṭe valayil oru srāvŭ peṭṭu.}
\end{quote}

\begin{quote}
A shark got caught in the fishermen's net.
\end{quote}

Fish markets on the coast sell small sharks too.

\begin{cousinweb}
\textbf{\ml{സ്രാവ്}} (\emph{srāvŭ}) is native Dravidian; DEDR lists the Malayalam word with Tamil \emph{cuṟā}, a shark (DEDR 2710).
\end{cousinweb}

\begin{grammarlens}[title={What you've built}]
One more word for getting around.
\end{grammarlens}

\subsection*{Your turn}
\begin{itemize}
\raggedright
  \item \emph{Say it:} \emph{srāvŭ}
  \item \emph{Say it:} \emph{srāvŭ}, once more
  \item \emph{Say it:} \emph{mīnpiṭuttakkāruṭe valayil oru srāvŭ peṭṭu.} — A shark got caught in the fishermen's net.
  \item \emph{Recall:} say the Malayalam for a swamp, then the Malayalam for a coral reef, then say \emph{srāvŭ} again
\end{itemize}

\begin{tcolorbox}[breakable,colback=teal!4,colframe=teal!35!black,title={Before you move on}]
What does \ml{സ്രാവ്} mean? (\textquotedblleft{}A shark\textquotedblright{}.) Say it once more.
\end{tcolorbox}
\section[\texorpdfstring{\ml{അടിഞ്ഞുകൂടുക} (aṭiññukūṭuka)}{അടിഞ്ഞുകൂടുക (aṭiññukūṭuka)}]{\ml{അടിഞ്ഞുകൂടുക} (aṭiññukūṭuka) — to accumulate; to settle (as silt)}

\label{lesson:ML-C498-atinnukutuka}



\begin{quote}
Before the new one: say the Malayalam for a coral reef, then the Malayalam for a shark.
\end{quote}

\subsection*{You'll want to know: \ml{അടിഞ്ഞുകൂടുക}}
\textbf{\ml{അടിഞ്ഞുകൂടുക}} — \emph{aṭiññukūṭuka} — \textquotedblleft{}to accumulate; to settle (as silt)\textquotedblright{}.

\begin{quote}
\ml{പുഴയുടെ} \ml{തീരത്ത്} \ml{ചെളി} \ml{അടിഞ്ഞുകൂടി}.
\end{quote}

\begin{quote}
\emph{puḻayuṭe tīrattŭ ceḷi aṭiññukūṭi.}
\end{quote}

\begin{quote}
Mud settled along the bank of the river.
\end{quote}

Past: \ml{അടിഞ്ഞുകൂടി}, \emph{aṭiññukūṭi}.

\ml{അടിയുക}, to settle at the bottom, plus \ml{കൂടുക}, to gather.

\begin{cousinweb}
\textbf{\ml{അടിഞ്ഞുകൂടുക}} (\emph{aṭiññukūṭuka}) is a native Dravidian compound built on \ml{അടി}, the bottom (DEDR 72), and \ml{കൂടുക}, to come together (DEDR 1882).
\end{cousinweb}

\begin{grammarlens}[title={What you've built}]
One more word for getting around.
\end{grammarlens}

\subsection*{Your turn}
\begin{itemize}
\raggedright
  \item \emph{Say it:} \emph{aṭiññukūṭuka}
  \item \emph{Say it:} \emph{aṭiññukūṭuka}, once more
  \item \emph{Say it:} \emph{puḻayuṭe tīrattŭ ceḷi aṭiññukūṭi.} — Mud settled along the bank of the river.
  \item \emph{Recall:} say the Malayalam for a coral reef, then the Malayalam for a shark, then say \emph{aṭiññukūṭuka} again
\end{itemize}

\begin{tcolorbox}[breakable,colback=teal!4,colframe=teal!35!black,title={Before you move on}]
What does \ml{അടിഞ്ഞുകൂടുക} mean? (\textquotedblleft{}To accumulate; to settle (as silt)\textquotedblright{}.) Say it once more.
\end{tcolorbox}
